% GENERATED FILE. Edit canonical lessons, then run npm run generate:books.
% canonical-source-hash: fnv1a64:bfaf028be12e4dd8
% canonical-lessons: SA-C93-kadapi, SA-C93-vilambah, SA-C93-vamah, SA-C93-daksinah, SA-C93-antah

\chapter{Ever, Delay, Left, Right, Inside}
\label{ch:sa-ever-delay-left-right-inside}
\hlchaptermodality{93}

\begin{chapteropening}
\textbf{By the end of this chapter:} \emph{I can say ever, a delay, left, right and inside.}

Complete the last lesson of chapter 93: i can say ever, a delay, left, right and inside.
\end{chapteropening}

\section[kadāpi]{\sk{कदापि} (kadāpi) — ever}

\label{lesson:SA-C93-kadapi}



\begin{quote}
Before the new one: say the Sanskrit for now, then the Sanskrit for always.
\end{quote}

\subsection*{You'll want to know: \sk{कदापि}}
\textbf{\sk{कदापि}} — \emph{kadāpi} — \textquotedblleft{}ever\textquotedblright{}.

\begin{grammarlens}[title={What you've built}]
One more word for this chapter.
\end{grammarlens}

\subsection*{Your turn}
\begin{itemize}
\raggedright
  \item \emph{Say it:} \emph{kadāpi}
  \item \emph{Say it:} \emph{kadāpi}, once more
  \item \emph{Say it:} say \emph{kadāpi}
  \item \emph{Recall:} say the Sanskrit for now, then the Sanskrit for always, then say \emph{kadāpi} again
\end{itemize}

\begin{tcolorbox}[breakable,colback=teal!4,colframe=teal!35!black,title={Before you move on}]
What does \sk{कदापि} mean? (\textquotedblleft{}Ever\textquotedblright{}.) Say it once more.
\end{tcolorbox}
\section[vilambaḥ]{\sk{विलम्बः} (vilambaḥ) — a delay}

\label{lesson:SA-C93-vilambah}



\begin{quote}
Before the new one: say the Sanskrit for always, then the Sanskrit for ever.
\end{quote}

\subsection*{You'll want to know: \sk{विलम्बः}}
\textbf{\sk{विलम्बः}} — \emph{vilambaḥ} — \textquotedblleft{}a delay\textquotedblright{}.

\begin{grammarlens}[title={What you've built}]
One more word for this chapter.
\end{grammarlens}

\subsection*{Your turn}
\begin{itemize}
\raggedright
  \item \emph{Say it:} \emph{vilambaḥ}
  \item \emph{Say it:} \emph{vilambaḥ}, once more
  \item \emph{Say it:} say \emph{vilambaḥ}
  \item \emph{Recall:} say the Sanskrit for always, then the Sanskrit for ever, then say \emph{vilambaḥ} again
\end{itemize}

\begin{tcolorbox}[breakable,colback=teal!4,colframe=teal!35!black,title={Before you move on}]
What does \sk{विलम्बः} mean? (\textquotedblleft{}A delay\textquotedblright{}.) Say it once more.
\end{tcolorbox}
\section[vāmaḥ]{\sk{वामः} (vāmaḥ) — left}

\label{lesson:SA-C93-vamah}



\begin{quote}
Before the new one: say the Sanskrit for ever, then the Sanskrit for a delay.
\end{quote}

\subsection*{You'll want to know: \sk{वामः}}
\textbf{\sk{वामः}} — \emph{vāmaḥ} — \textquotedblleft{}left\textquotedblright{}.

\begin{grammarlens}[title={What you've built}]
One more word for this chapter.
\end{grammarlens}

\subsection*{Your turn}
\begin{itemize}
\raggedright
  \item \emph{Say it:} \emph{vāmaḥ}
  \item \emph{Say it:} \emph{vāmaḥ}, once more
  \item \emph{Say it:} say \emph{vāmaḥ}
  \item \emph{Recall:} say the Sanskrit for ever, then the Sanskrit for a delay, then say \emph{vāmaḥ} again
\end{itemize}

\begin{tcolorbox}[breakable,colback=teal!4,colframe=teal!35!black,title={Before you move on}]
What does \sk{वामः} mean? (\textquotedblleft{}Left\textquotedblright{}.) Say it once more.
\end{tcolorbox}
\section[dakṣiṇaḥ]{\sk{दक्षिणः} (dakṣiṇaḥ) — right; south}

\label{lesson:SA-C93-daksinah}



\begin{quote}
Before the new one: say the Sanskrit for a delay, then the Sanskrit for left.
\end{quote}

\subsection*{You'll want to know: \sk{दक्षिणः}}
\textbf{\sk{दक्षिणः}} — \emph{dakṣiṇaḥ} — \textquotedblleft{}right; south\textquotedblright{}.

\begin{grammarlens}[title={What you've built}]
One more word for this chapter.
\end{grammarlens}

\subsection*{Your turn}
\begin{itemize}
\raggedright
  \item \emph{Say it:} \emph{dakṣiṇaḥ}
  \item \emph{Say it:} \emph{dakṣiṇaḥ}, once more
  \item \emph{Say it:} say \emph{dakṣiṇaḥ}
  \item \emph{Recall:} say the Sanskrit for a delay, then the Sanskrit for left, then say \emph{dakṣiṇaḥ} again
\end{itemize}

\begin{tcolorbox}[breakable,colback=teal!4,colframe=teal!35!black,title={Before you move on}]
What does \sk{दक्षिणः} mean? (\textquotedblleft{}Right; south\textquotedblright{}.) Say it once more.
\end{tcolorbox}
\section[antaḥ]{\sk{अन्तः} (antaḥ) — inside}

\label{lesson:SA-C93-antah}



\begin{quote}
Before the new one: say the Sanskrit for left, then the Sanskrit for right; south.
\end{quote}

\subsection*{You'll want to know: \sk{अन्तः}}
\textbf{\sk{अन्तः}} — \emph{antaḥ} — \textquotedblleft{}inside\textquotedblright{}.

\begin{grammarlens}[title={What you've built}]
One more word for this chapter.
\end{grammarlens}

\subsection*{Your turn}
\begin{itemize}
\raggedright
  \item \emph{Say it:} \emph{antaḥ}
  \item \emph{Say it:} \emph{antaḥ}, once more
  \item \emph{Say it:} say \emph{antaḥ}
  \item \emph{Recall:} say the Sanskrit for left, then the Sanskrit for right; south, then say \emph{antaḥ} again
\end{itemize}

\begin{tcolorbox}[breakable,colback=teal!4,colframe=teal!35!black,title={Before you move on}]
What does \sk{अन्तः} mean? (\textquotedblleft{}Inside\textquotedblright{}.) Say it once more.
\end{tcolorbox}
